%% -------------------------------------------------------
%% Kapitel 6: iOS-Applikation
%% -------------------------------------------------------
\chapter{iOS-Applikation}
\label{chap:ios}

\section{Architektur der iOS-App}
\label{sec:ios-architektur}

Die iOS-Applikation wurde mit \textbf{Swift} und \textbf{SwiftUI} entwickelt und folgt
dem \textbf{MVVM-Muster} (Model-View-ViewModel). Diese Architektur trennt die
UI-Darstellung (View) von der Geschäftslogik und dem Zustand (ViewModel) und erleichtert
Testbarkeit und Wartbarkeit.

\begin{description}
  \item[Model] Entspricht den Datenstrukturen, die der Backend-API entsprechen
    (\texttt{Dog}, \texttt{Owner}, \texttt{Breed}, \texttt{LegalRestrictions}).
  \item[ViewModel] Verwaltet den Zustand der UI mit \texttt{@Published}-Eigenschaften
    und kommuniziert mit dem Backend über HTTP. Durch das \texttt{@StateObject}-Macro
    wird der Lebenszyklus an die View gebunden.
  \item[View] Deklarative SwiftUI-Views lesen den Zustand aus dem ViewModel und reagieren
    automatisch auf Änderungen.
\end{description}

\begin{figure}[H]
  \centering
  \begin{tikzpicture}[
      box/.style={draw, rectangle, rounded corners=4pt,
                  minimum width=4.0cm, minimum height=0.9cm,
                  align=center, font=\small, fill=gray!10, text width=3.8cm},
      ext/.style={draw, rectangle, rounded corners=4pt,
                  minimum width=4.0cm, minimum height=0.9cm,
                  align=center, font=\small, fill=white, dashed, text width=3.8cm},
      arr/.style={{Stealth[length=5pt]}-{Stealth[length=5pt]}, thick},
      oarr/.style={-{Stealth[length=5pt]}, thick, dashed}
    ]
    \node[box] (view)  at (0, 0)   {View\\(SwiftUI)};
    \node[box] (vm)    at (6, 0)   {ViewModel\\(@ObservableObject)};
    \node[box] (model) at (12, 0)  {Model\\(Codable DTOs)};
    \node[ext] (api)   at (6,-2)   {Backend REST-API\\(URLSession)};

    \draw[arr]  (view)  -- node[font=\footnotesize, above]{Bindings} (vm);
    \draw[arr]  (vm)    -- node[font=\footnotesize, above]{Daten}    (model);
    \draw[oarr] (vm)    -- node[font=\footnotesize, right]{HTTP}     (api);
  \end{tikzpicture}
  \caption{MVVM-Architektur der DoggoGO iOS-App}
  \label{fig:mvvm}
\end{figure}

\section{Projektstruktur}
\label{sec:ios-struktur}

Das Xcode-Projekt ist in folgende Gruppen gegliedert:

\begin{itemize}
  \item \texttt{Models/} – \texttt{Codable}-Strukturen für API-DTOs
  \item \texttt{ViewModels/} – \texttt{ObservableObject}-Klassen mit Geschäftslogik
  \item \texttt{Views/} – SwiftUI-Views (Login, Map, DogList, BreedPicker, etc.)
  \item \texttt{Services/} – \texttt{APIService} (HTTP-Kommunikation), \texttt{AuthService}
    (Keychain-Zugriff)
  \item \texttt{Assets.xcassets} – App-Icons, Farbschemata
\end{itemize}

\section{MapKit-Integration}
\label{sec:mapkit}

Die Kartenansicht basiert auf \textbf{MapKit}, Apples nativer Kartenbibliothek. In
SwiftUI wird eine Karte mit \texttt{Map(coordinateRegion:)} eingebunden. Der aktuelle
Standort der Nutzerin / des Nutzers wird über \texttt{CLLocationManager} ermittelt und
als Region übergeben.

\begin{lstlisting}[language=swift, caption={Kartenansicht in SwiftUI (vereinfacht)}, label=lst:mapview, float=h]
import SwiftUI
import MapKit

struct MapView: View {
    @StateObject private var locationManager = LocationManager()
    @State private var region = MKCoordinateRegion(
        center: CLLocationCoordinate2D(
            latitude: 48.2692, longitude: 14.2956), // Leonding
        span: MKCoordinateSpan(
            latitudeDelta: 0.05, longitudeDelta: 0.05)
    )

    var body: some View {
        Map(coordinateRegion: $region,
            showsUserLocation: true,
            annotationItems: locationManager.pois) { poi in
            MapMarker(coordinate: poi.coordinate,
                      tint: poi.color)
        }
        .ignoresSafeArea()
        .onReceive(locationManager.$userLocation) { loc in
            guard let loc else { return }
            region.center = loc.coordinate
        }
    }
}
\end{lstlisting}

Für die Darstellung von Points of Interest (Freilaufwiesen, Sackerlspender,
Hundekot-Mülleimer) werden \texttt{MapMarker}-Annotationen mit farblicher Codierung
pro Kategorie genutzt.

\section{Standortdienste und Berechtigungen}
\label{sec:standortdienste}

Die Standortermittlung wird über \texttt{CLLocationManager} realisiert. Gemäß
DSGVO-Anforderung (NF-03) darf der Standort nur mit expliziter Nutzer\-innen-/Nut\-zer\-zu\-stim\-mung
erfasst werden. Die App fragt beim ersten Start die Berechtigung
\texttt{re\-quest\-When\-In\-Use\-Au\-thor\-iza\-tion()} ab. Die \texttt{Info.plist} enthält die
Pflichteinträge \texttt{NS\-Lo\-ca\-tion\-When\-In\-Use\-Usage\-De\-scrip\-tion} mit einer verständlichen
Begründung.

\begin{lstlisting}[language=swift, caption={LocationManager (Auszug)}, label=lst:locationmanager, float=h]
class LocationManager: NSObject, ObservableObject, CLLocationManagerDelegate {
    private let manager = CLLocationManager()
    @Published var userLocation: CLLocation?

    override init() {
        super.init()
        manager.delegate = self
        manager.desiredAccuracy = kCLLocationAccuracyBest
        manager.requestWhenInUseAuthorization()
        manager.startUpdatingLocation()
    }

    func locationManager(_ manager: CLLocationManager,
                         didUpdateLocations locations: [CLLocation]) {
        userLocation = locations.last
    }
}
\end{lstlisting}

\section{Authentifizierung in der iOS-App}
\label{sec:ios-auth}

\subsection{Login und Registrierung}
\label{subsec:ios-login}

Die Anmeldeansicht (\texttt{LoginView}) enthält Eingabefelder für E-Mail und Passwort.
Nach dem Tippen auf „Anmelden" wird \texttt{POST /api/owners/login} mit den Credentials
aufgerufen. Bei Erfolg liefert das Backend ein JWT zurück:

\begin{lstlisting}[language=swift, caption={Login-Anfrage (Auszug)}, label=lst:login, float=h]
func login(email: String, password: String) async throws {
    let body = LoginRequest(email: email, password: password)
    let response: AuthResponse = try await apiService.post(
        endpoint: "/api/owners/login", body: body)
    AuthService.shared.saveToken(response.token)
    isAuthenticated = true
}
\end{lstlisting}

\subsection{Token-Persistenz im Keychain}
\label{subsec:keychain}

Das JWT wird sicher im iOS~\texttt{Keychain} gespeichert, nicht in
\texttt{UserDefaults}, da der Keychain verschlüsselt und vor anderen Apps geschützt ist.
Der \texttt{AuthService} kapselt alle Keychain-Operationen:

\begin{lstlisting}[language=swift, caption={AuthService – Keychain-Speicherung (Auszug)}, label=lst:keychain, float=h]
class AuthService {
    static let shared = AuthService()
    private let tokenKey = "doggo_jwt_token"

    func saveToken(_ token: String) {
        let data = Data(token.utf8)
        let query: [String: Any] = [
            kSecClass as String: kSecClassGenericPassword,
            kSecAttrAccount as String: tokenKey,
            kSecValueData as String: data
        ]
        SecItemDelete(query as CFDictionary)
        SecItemAdd(query as CFDictionary, nil)
    }

    func getToken() -> String? {
        let query: [String: Any] = [
            kSecClass as String: kSecClassGenericPassword,
            kSecAttrAccount as String: tokenKey,
            kSecReturnData as String: true
        ]
        var item: CFTypeRef?
        guard SecItemCopyMatching(query as CFDictionary, &item) == noErr,
              let data = item as? Data else { return nil }
        return String(data: data, encoding: .utf8)
    }
}
\end{lstlisting}

Beim App-Start prüft das Root-ViewModel, ob ein Token im Keychain vorhanden ist, und
überspringt in diesem Fall die Anmeldeseite direkt zur Kartenansicht.

\section{Hundeprofilverwaltung}
\label{sec:hundeprofil}

Die \texttt{DogListView} listet alle Hunde der angemeldeten Nutzerin / des angemeldeten
Nutzers auf und ruft dazu \texttt{GET /api/dogs} mit dem Bearer-Token auf. Über einen
\texttt{NavigationLink} gelangt man zur Detailansicht oder zum Formular zum Anlegen /
Bearbeiten eines Profils. Die Rassenauswahl erfolgt über einen \texttt{Picker}, der die
von \texttt{GET /api/breeds} geladene Liste anzeigt.

\begin{lstlisting}[language=swift, caption={DogListView – Übersicht der Hundeprofile (Auszug)}, label=lst:doglist, float=h]
struct DogListView: View {
    @StateObject private var viewModel = DogListViewModel()

    var body: some View {
        NavigationView {
            List(viewModel.dogs) { dog in
                NavigationLink(destination: DogDetailView(dog: dog)) {
                    VStack(alignment: .leading) {
                        Text(dog.name).font(.headline)
                        Text("\(dog.breed.breedName), \(dog.age) Jahre")
                            .font(.subheadline).foregroundColor(.secondary)
                    }
                }
            }
            .navigationTitle("Meine Hunde")
            .toolbar {
                ToolbarItem(placement: .navigationBarTrailing) {
                    NavigationLink("Hinzufuegen",
                        destination: AddDogView())
                }
            }
            .task { await viewModel.loadDogs() }
        }
    }
}
\end{lstlisting}

\section{Anzeige gesetzlicher Vorschriften}
\label{sec:ios-legal}

Nach der Rassenauswahl ruft die App \texttt{GET /api/breeds/\{id\}} ab und zeigt die
\texttt{LegalRestrictions} in einer strukturierten Ansicht an. Rassen der Kategorie
\texttt{spezi\-elle\_hun\-de\-rasse} werden farblich hervorgehoben (rotes Label), um die
Nutzerin / den Nutzer auf besondere Pflichten aufmerksam zu machen:

\begin{lstlisting}[language=swift, caption={LegalRestrictionsView – Anzeige der Vorschriften}, label=lst:legal, float=h]
struct LegalRestrictionsView: View {
    let restrictions: LegalRestrictions

    var body: some View {
        VStack(alignment: .leading, spacing: 8) {
            Text("Rechtliche Vorgaben")
                .font(.title2).bold()
            if restrictions.category == "spezielle_hunderasse" {
                Text("Spezielle Hunderasse!")
                    .foregroundColor(.red).bold()
            }
            ForEach(restrictions.rules, id: \.self) { rule in
                Label(rule, systemImage: "exclamationmark.circle")
                    .font(.callout)
            }
        }
        .padding()
    }
}
\end{lstlisting}

\section{HTTP-Kommunikation mit dem Backend}
\label{sec:ios-http}

Alle Netzwerkanfragen werden über einen zentralen \texttt{APIService} abgewickelt, der
\texttt{URLSession} kapselt. Das Bearer-Token wird in jedem Request automatisch als
\texttt{Authorization}-Header eingefügt:

\begin{lstlisting}[language=swift, caption={APIService – Authentifizierter Request (Auszug)}, label=lst:apiservice, float=h]
func get<T: Decodable>(endpoint: String) async throws -> T {
    guard let url = URL(string: baseURL + endpoint) else {
        throw APIError.invalidURL
    }
    var request = URLRequest(url: url)
    if let token = AuthService.shared.getToken() {
        request.setValue("Bearer \(token)",
                         forHTTPHeaderField: "Authorization")
    }
    let (data, response) = try await URLSession.shared.data(for: request)
    guard let http = response as? HTTPURLResponse,
          (200..<300).contains(http.statusCode) else {
        throw APIError.serverError
    }
    return try JSONDecoder().decode(T.self, from: data)
}
\end{lstlisting}

\section{Fehlerbehandlung und Ladezustände}
\label{sec:ios-fehler}

Eine robuste iOS-App muss mit Netzwerkfehlern, ungültigen Tokens und unerwarteten
Serverantworten umgehen. In der DoggoGO-App wird Fehlerbehandlung konsequent im
ViewModel-Layer zentralisiert – die View selbst kennt keine Netzwerklogik.

\subsection{Fehlerdefinition}
\label{subsec:api-error}

Alle möglichen Fehlerszenarien werden in einem \texttt{APIError}-Enum typisiert:

\begin{lstlisting}[language=swift, caption={APIError – Typisierte Fehlerdefinition}, label=lst:apierror]
enum APIError: LocalizedError {
    case invalidURL
    case unauthorized       // 401 – JWT abgelaufen oder ungueltig
    case notFound           // 404 – Ressource nicht vorhanden
    case serverError        // 5xx – Backend-Fehler
    case decodingError      // JSON-Parsing fehlgeschlagen
    case noConnection       // Netzwerk nicht erreichbar

    var errorDescription: String? {
        switch self {
        case .invalidURL:    return "Ungueltige URL."
        case .unauthorized:  return "Sitzung abgelaufen. Bitte erneut anmelden."
        case .notFound:      return "Ressource nicht gefunden."
        case .serverError:   return "Serverfehler. Bitte spaeter versuchen."
        case .decodingError: return "Antwort konnte nicht verarbeitet werden."
        case .noConnection:  return "Keine Internetverbindung."
        }
    }
}
\end{lstlisting}

Das Protokoll \texttt{LocalizedError} stellt sicher, dass SwiftUI und das Betriebssystem
die Fehlerbeschreibung automatisch für Alert-Dialoge verwenden können.

\subsection{Ladezustand und Fehleranzeige im ViewModel}
\label{subsec:loading-state}

Jedes ViewModel hält zwei Zustandsvariablen: \texttt{isLoading} blendet während einer
laufenden Netzwerkanfrage eine Ladeindikator-Ansicht ein, \texttt{errorMessage} speichert
eine lesbare Fehlerbeschreibung für den Alert-Dialog:

\begin{lstlisting}[language=swift, caption={ViewModel – Ladezustand und Fehlerbehandlung}, label=lst:loading]
class DogListViewModel: ObservableObject {
    @Published var dogs: [Dog] = []
    @Published var isLoading = false
    @Published var errorMessage: String?

    func loadDogs() async {
        isLoading = true
        defer { isLoading = false }
        do {
            dogs = try await APIService.shared.get(endpoint: "/api/dogs")
        } catch APIError.unauthorized {
            errorMessage = APIError.unauthorized.errorDescription
            AuthService.shared.deleteToken()   // abgelaufenes Token loeschen
        } catch {
            errorMessage = error.localizedDescription
        }
    }
}
\end{lstlisting}

Das \texttt{defer}-Statement stellt sicher, dass \texttt{isLoading} immer auf
\texttt{false} zurückgesetzt wird – auch wenn die Anfrage mit einem Fehler endet.
Bei einem \texttt{401~Unauthorized} wird das gespeicherte JWT automatisch gelöscht,
sodass die Nutzerin / der Nutzer zur Anmeldeseite weitergeleitet wird.

\subsection{Alert-Anzeige in der View}
\label{subsec:alert}

Die View abonniert \texttt{errorMessage} über ein \texttt{.alert}-Modifier und zeigt
bei gesetztem Wert automatisch einen Dialog an:

\begin{lstlisting}[language=swift, caption={SwiftUI – Alert fuer Fehlermeldungen}, label=lst:alert]
.alert("Fehler", isPresented: Binding(
    get: { viewModel.errorMessage != nil },
    set: { if !$0 { viewModel.errorMessage = nil } }
)) {
    Button("OK", role: .cancel) { }
} message: {
    Text(viewModel.errorMessage ?? "")
}
.overlay {
    if viewModel.isLoading {
        ProgressView("Laden ...")
            .padding()
            .background(.regularMaterial, in: RoundedRectangle(cornerRadius: 8))
    }
}
\end{lstlisting}

Dieser Ansatz hält die View vollständig frei von Fehlerlogik: Sie bindet lediglich den
Zustand des ViewModels und delegiert jede Entscheidung dorthin.
